\section{Ricci growth at the recognition endpoint}\label{sec:endpoint-windows}

Let \(I=T=4\pi^2\) and retain the fixed physical minimizing
sequence of this paper. Smooth recognition implies that every
\(d_i=e_i-T\) is positive, since a smooth positive class of energy
at most \(T\) would give \(I=0\). The strict defect inequality
\eqref{eq:endpoint-deficit}, or independently the full-tensor
pinching classification of Fu \cite[Theorem 1.3]{Fu2018FourManifolds},
together with \eqref{eq:budgets} gives, with \(A_i=\|Z_{g_i}\|_{L^2(g_i)}\),
\begin{equation}\label{eq:endpoint-ricci-defect}
 A_i^2\le 8d_i,\qquad
 192\pi^2-48d_i\le Y_i^2\le192\pi^2+48d_i .
\end{equation}
In particular the physical trace-free Ricci energy tends to zero.

For a physical region \(U\), let
\[
 N_{i,q}(U)=\sup_{t>0}t\,\operatorname{Vol}_{g_i}
      (U\cap\{|Z_{g_i}|_{g_i}>t\})^{1/q},\qquad
 \gamma_q=2-4/q .
\]
This is the weak \(L^q\) quasinorm in the physical gauge.

\begin{proposition}[Ricci growth and curvature at the endpoint]
\label{prop:endpoint-ricci-windows}
For every fixed finite \(q>2\) and \(\kappa>0\), a physical
minimizing sequence at \(I=T\) satisfies
\[
 A_i^{2\kappa}N_{i,q}(M)\longrightarrow\infty,\qquad
 d_i^\kappa N_{i,q}(M)\longrightarrow\infty.
\]
In particular \(d_i^\kappa\sup_M|\Rm_{g_i}|_{\rm tens}\to\infty\).
Let \(a_i\to0\), \(x_i\in M\), and
\(U_i(D)=B_{g_i}(x_i,Da_i)\). If for every fixed finite \(D>0\)
there is a bound, eventually on one subsequence, on
\begin{equation}\label{eq:endpoint-protected-weak}
 \|Z_{g_i}\|_{L^2(U_i(D),g_i)}^{2\kappa}
       a_i^{\gamma_q}N_{i,q}(U_i(D)),
\end{equation}
then for each fixed \(A>0\), the full curvature integral on
\(B_{g_i}(x_i,Aa_i)\) tends to zero on that subsequence. One
sufficient condition for \eqref{eq:endpoint-protected-weak} is a
bound with \(d_i^\kappa\) in place of the factor involving the
local \(L^2\) norm.

For \(a_i=(d_i^\kappa N_{i,q}(M))^{-1/\gamma_q}\) the curvature
conclusion holds uniformly in the center. If \(r_i(\vartheta)\)
is the least physical radius at which some ball carries positive
chiral mass \(\vartheta e_i\), then
\[
 r_i(\vartheta)(d_i^\kappa N_{i,q}(M))^{1/\gamma_q}
     \longrightarrow\infty\qquad (0<\vartheta\le1/2).
\]
\end{proposition}

\begin{proof}
On any measure space put \(A=\|f\|_2\), \(B=\|f\|_{q,\infty}\).
The distribution function of \(|f|\) is bounded by
\(\min(A^2t^{-2},B^qt^{-q})\). For \(2<p<q\), split its
layer-cake integral at \(t=(B^q/A^2)^{1/(q-2)}\), obtaining
\[
 \|f\|_p\le C(p,q)A^\alpha B^{1-\alpha},\qquad
 \frac1p=\frac{\alpha}{2}+\frac{1-\alpha}{q} .
\]
If either norm vanishes both sides are zero. Choose
\(p=2+(q-2)/(1+q\kappa)\), so that
\(\alpha/(1-\alpha)=2\kappa\) and
\(\gamma_p=(1-\alpha)\gamma_q\). Therefore
\[
 a_i^{\gamma_p}\|Z_{g_i}\|_{L^p(U,g_i)}
 \le C(p,q)\left(\|Z_{g_i}\|_{L^2(U,g_i)}^{2\kappa}
          a_i^{\gamma_q}N_{i,q}(U)\right)^{1-\alpha} .
\]
On \(M\), a bounded subsequence of \(A_i^{2\kappa}N_{i,q}(M)\)
would imply a bounded subsequence of the global \(L^p\) norm.
The physical Ricci-scale theorem
gives smooth positive attainment on the original type, which
contradicts recognition at \(T\). The bound
\eqref{eq:endpoint-ricci-defect} gives the stated divergence in
terms of \(d_i\). Unit volume and the pointwise
Ricci contraction bound give \(N_{i,q}(M)\le C\sup_M|\Rm|_{\rm tens}\).

On the shrinking balls, the displayed interpolation gives the
strong \(L^p\) bounds on all protected windows required by the
physical Ricci-scale theorem in the companion master manuscript.
Its conclusion is vanishing full curvature mass on every smaller
fixed multiple of the scale. For the specified \(a_i\), the
sufficient bound involving \(d_i^\kappa\) holds uniformly over
the centers for every \(D\). Failure of uniform vanishing would
select a subsequence and centers with a positive integral. If a
bounded multiple of \(a_i\) reached a chiral quantile, its full
curvature integral would be at least \(4\vartheta e_i\), contrary
to this vanishing.
\end{proof}

This condition allows growth of the weak \(L^q\) norm controlled
by any fixed inverse power of the minimizing excess. The
critical estimate \eqref{eq:endpoint-ricci-defect} itself gives
neither an exponent above two nor the coordinate regularity
needed by the harmonic compactness argument. A candidate endpoint
must meet that critical estimate while its supercritical norms
grow faster than every such fixed power.
